%% ++++++++++++++++++++++++++++++++++++++++++++++++++++++++++++
%% Hauptdatei, Wurzel des Dokuments
%% ++++++++++++++++++++++++++++++++++++++++++++++++++++++++++++
%
%  Gerüst:
%  * Version 0.13
%  * M.Sc. Silvia Krug, silvia.krug@tu-ilmenau.de
%  * Fachgebiet Kommunikationsnetze, TU Ilmenau
%
%  Für Hauptseminare, Studienarbeiten, Diplomarbeiten
%
%  Autor           : Max Mustermann
%  Letzte Änderung : 04.12.2015
%

% Headerfeld, Typ des Dokumentes, einzubindende Packages.
% Hier bei Bedarf Änderungen vornehmen.
\documentclass
[   twoside=false,     % Einseitiger oder zweiseitiger Druck?
    fontsize=12pt,     % Bezug: 12-Punkt Schriftgröße
    DIV=15,            % Randaufteilung, siehe Dokumentation "KOMA"-Script
    BCOR=17mm,         % Bindekorrektur: Innen 17mm Platz lassen. Copyshop-getestet.
%    headsepline,
    headsepline,  % Unter Kopfzeile Trennlinie (aus: headnosepline)
    footsepline,  % Über Fußzeile Trennlinie (aus: footnosepline)
    open=right,        % Neue Kapitel im zweiseitigen Druck rechts beginnen lassen
    paper=a4,          % Seitenformat A4
    abstract=true,     % Abstract einbinden
    listof=totoc,      % Div. Verzeichnisse ins Inhaltsverzeichnis aufnehmen
    bibliography=totoc,% Literaturverzeichnis ins Inhaltsverzeichnis aufnehmen
    titlepage,         % Titelseite aktivieren
    headinclude=true,  % Seiten-Head in die Satzspiegelberechnung mit einbeziehen
    footinclude=false, % Seiten-Foot nicht in die Satzspiegelberechnung mit einbeziehen
    numbers=noenddot   % Gliederungsnummern ohne abschließenden Punkt darstellen
]   {scrreprt}         % Dokumentenstil: "Report" aus dem KOMA-Skript-Paket

\usepackage[active]{srcltx}
\usepackage{ifthen}
\usepackage[utf8]{inputenc}  
\usepackage[T1]{fontenc}
\usepackage{lmodern}
\usepackage[T1]{url}
\usepackage[final]{graphicx}
\usepackage[automark]{scrpage2}
\usepackage{setspace}
\usepackage[plainpages=false,pdfpagelabels,hypertexnames=false]{hyperref}
\usepackage{amsmath,amssymb,amsfonts}  
\usepackage{bm}                        
\usepackage{graphicx}                   
\usepackage{psfrag}                   
\usepackage{array}
\usepackage{caption}
\usepackage{hyperref}
\usepackage{stfloats}
\usepackage{textcomp}
\usepackage{tabulary}
\usepackage{amsfonts}
\usepackage{epstopdf}
\usepackage{float}


\newcommand{\bs}{\boldsymbol}
\newcommand{\compl}{\mathbb{C}}      
\newcommand{\real}{\mathbb{R}} 
\newcommand{\artderausarbeitung}{Masterarbeit}
\newcommand{\namedesautors}{Bulat Valeev}
\newcommand{\themaderarbeit}{Tensor based system modelling for GFDM}


% Tiefe der Kapitelnummerierung beeinflussen
\setcounter{secnumdepth}{3} % Tiefe der Nummerierung
\setcounter{tocdepth}{3}    % Tiefe des Inhaltsverzeichnisses
% PDF Metadaten definieren
\hypersetup{
   pdftitle={\themaderarbeit},
   pdfsubject={\artderausarbeitung},
   pdfauthor={\namedesautors},
   pdfkeywords={\artderausarbeitung; TU-Ilmenau; }}


\usepackage[intoc]{nomencl}
  \AtBeginDocument{\setlength{\nomlabelwidth}{.25\columnwidth}}
  \let\abbrev\nomenclature
  \renewcommand{\nomname}{Acronyms}
  \renewcommand{\nomlabel}[1]{#1 \dotfill}
  \setlength{\nomitemsep}{-\parsep}
  \makenomenclature

\usepackage[normalem]{ulem}
  \newcommand{\markup}[1]{\textbf{#1}}

% Seitenlayout festlegen. Hier nichts ändern!
\pagestyle{scrplain}
\ihead[]{\headmark}
\ohead[]{\pagemark}
\chead[]{}
\ifoot[]{}
\ofoot[]{\scriptsize \artderausarbeitung\ \namedesautors}
\cfoot[]{}
\renewcommand{\titlepagestyle}{scrheadings}
\renewcommand{\partpagestyle}{scrheadings}
\renewcommand{\chapterpagestyle}{scrheadings}
\renewcommand{\indexpagestyle}{scrheadings}

% Abschnittsweise Nummerierung anstatt fortlaufend. Hier nichts ändern!
\makeatletter
\@addtoreset{equation}{chapter}
\@addtoreset{figure}{chapter}
\@addtoreset{table}{chapter}
\renewcommand\theequation{\thechapter.\@arabic\c@equation}
\renewcommand\thefigure{\thechapter.\@arabic\c@figure}
\renewcommand\thetable{\thechapter.\@arabic\c@table}
\makeatother

% Quelltextrahmen, klein. Hier nichts ändern!
\newsavebox{\inhaltkl}
\def\rahmenkl{\sbox{\inhaltkl}\bgroup\small\renewcommand{\baselinestretch}{1}\vbox\bgroup\hsize\textwidth}
\def\endrahmenkl{\par\vskip-\lastskip\egroup\egroup\fboxsep3mm%
\framebox[\textwidth][l]{\usebox{\inhaltkl}}}

% Quelltextrahmen, normale Groesse. Hier nichts ändern!
\newsavebox{\inhalt}
\def\rahmen{\sbox{\inhalt}\bgroup\renewcommand{\baselinestretch}{1}\vbox\bgroup\hsize\textwidth}
\def\endrahmen{\par\vskip-\lastskip\egroup\egroup\fboxsep3mm%
\framebox[\textwidth][l]{\usebox{\inhalt}}}

% Trennvorschläge für falsch getrennte Wörter.
% Wird häufig bei eingedeutschen Wörtern benötigt, da LaTeX hierbei
% gerne falsch trennt. Alternativ kann auch im Fliesstext ein
% Trennvorschlag per "\-" hinterlegt werden, bspw.:
% Die Hard\-ware besteht aus A und B.
\hyphenation{
Hard-ware
}

% Sonstige Befehlsdefinitionen hier ablegen.
\newcommand{\entspricht}{\stackrel{\wedge}{=}}

% Tabellenspaltendefinitionen mit fester Breite --> somit Zeilenumbruch innerhalb einer Zelle möglich
% aus http://www.torsten-schuetze.de/tex/tabsatz-2004.pdf
\usepackage{array, booktabs}
\newcolumntype{f}{>{$}l<{$}}
\newcolumntype{n}{>{\raggedright}l}
\newcolumntype{N}{>{\scriptsize}l}
\newcolumntype{v}[1]{>{\raggedright\hspace{0pt}}m{#1}}
\newcolumntype{V}[1]{>{\scriptsize\raggedright\hspace{0pt}}m{#1}}
\newcolumntype{Z}[1]{>{\raggedright\centering}m{#1}}
\newcolumntype{k}[1]{>{\raggedright}p{#1}}
% ergibt Tabllenspalte fester Breite, linksbündig
% Umbruch innerhalb der Zelle mit \\, neue Tabellezeile mit \tabularnewline
% \addlinespace für Gruppentrennung (aus \texttt{booktabs.sty})


\begin{document}
\onehalfspacing
%%% ++++++++++++++++++++++++++++++++++++++++++++++++++++++++++++
%% Hauptdatei, Wurzel des Dokuments
%% ++++++++++++++++++++++++++++++++++++++++++++++++++++++++++++
%
%  Gerüst:
%  * Version 0.13
%  * M. Sc. Silvia Krug, silvia.krug@tu-ilmenau.de
%  * Fachgebiet Kommunikationsnetze, TU Ilmenau
%
%  Für Hauptseminare, Studienarbeiten, Diplomarbeiten
%
%  Autor           : Bulat Valeev
%  Letzte Änderung : 04.12.2015
%

\begin{titlepage}
\centering
\includegraphics[scale=0.5]{bilder/tui_logo}\\[3ex]
{\Large \textsc{Technische Universität Ilmenau}}\\[3ex]
{\Large Fakultät für Elektrotechnik und Informationstechnik}\\[3ex]
\vfill
{\Large \textbf{\artderausarbeitung}}\\[4ex]
{\large \textbf{\themaderarbeit}}\\[5ex]
\vfill
\begin{tabular}{rl}
\hline\\
vorgelegt von:          & \quad \namedesautors\\[1,5ex]
%eingereicht am:         & \quad 31.\,12.\,2011\\[1,5ex]
geboren am:             & \quad 13.\,06.\,1991 in Kazan\\[1,5ex]
Studiengang:            & \quad Elektrotechnik und Informationstechnik\\[1,5ex]
Studienrichtung:        & \quad Communication and signal processing \\[1,5ex]
%                        & \quad Kommunikationssysteme\\[5ex]
Anfertigung im Fachgebiet:
                        & \quad Communication and signal processing\\[1,5ex]
                        & \quad Fakultät für Elektrotechnik und Informationstechnik\\[1,5ex]
Verantwortlicher Professor:
                        & \quad Univ.-Prof.~Dr.-ing.Martin Haardt\\[1,5ex]
Wissenschaftlicher Betreuer:
                        & \quad M.~Sc. Kristina Naskovska\\
Datum:                       
						& \quad 15 Aug 2016\\[0,5ex]
\end{tabular}
\vfill
\end{titlepage}








%%% ++++++++++++++++++++++++++++++++++++++++++++++++++++++++++++
%% Hauptdatei, Wurzel des Dokuments
%% ++++++++++++++++++++++++++++++++++++++++++++++++++++++++++++
%
%  Gerüst:
%  * Version 0.13
%  * M.Sc. Silvia Krug, silvia.krug@tu-ilmenau.de
%  * Fachgebiet Kommunikationsnetze, TU Ilmenau
%
%  Für Hauptseminare, Studienarbeiten, Diplomarbeiten
%
%  Autor           : Max Mustermann
%  Letzte Änderung : 04.12.2015
%

\section*{Acknolegment}
 \emph{Firstly, I would like to express my sincere gratitude to the GRIAT organization personal from the Kazan. There are a lot of them, thanks for giving me chance to start this thesis\linebreak 
Secondly, I am very grateful to Kristina Naskovska for help and supervision. Without her advice, most of this work would not have been possible\linebreak 
Also I am very grateful to the "RIIT" department in the University where I studied the bachelor degree. They supported a lot throughout my master research.\linebreak 
Finally I would like to thank my family for constantly supporting me in work during the last half of year in Kazan. It is the great success for me to achieve this point in the life}


\renewcommand{\abstractname}{Abstract}
\begin{abstract}
\emph{The radio resourses bacome more and more expensive nowadays. The requirements for the transmission systems become more and more demanding. The necessary transmission rate also become more and more high. This demand increase complexity for a designed systems. \\
Generalized frequency division modulation is the new approach to increase spectral efficiency in comparison to the modern OFDM scheme.  GFDM uses time mixing for the symbols in one transmission block to increase density of the transmitted data. The new approach is flexible and can adjust a frequency overlap with one variable. The frequency overlap slightly increase the SER of the system, but decrease out-of-band radiation. The system also can adjust the number of turned on subcarriers in case if frequency sharing used. The system can wok in the same frequency bands as ther legal users and doesn't interfere with them. The transmission system become complex in the symbol estimation, but introduce the self-interference between symbols.\\
GFDM modulation can be explained with the PARATUCK2 tensor model to present modulation scheme in more efficient and natural way. This model explain the transmitted data as the slice-wise product between three different matrices and two tensors. This notation allow to introduce time and frequency domain separately for the symbols in one transmission block. The PARATUCK2 model can be exploited to find more efficient ways to estimate unknown symbols or update channel variables}
\end{abstract}
%
\cleardoublepage
\ihead[]{List of Figures}
\listoffigures

%
\cleardoublepage
\ihead[]{List of Tables}
\listoftables


%\cleardoublepage % Seitenumbruch erzwingen vor Änderung des Nummerierungsstils
%\pagenumbering{roman} % Nummerierung der Seiten ab hier: i, ii, iii, iv...
%\pagestyle{scrheadings} % Ab hier mit Kopf- und Fusszeile
%\tableofcontents

% Die einzelnen Kapitel
%\cleardoublepage % Seitenumbruch erzwingen vor Änderung des Nummerierungsstils
%\pagenumbering{arabic} % Nummerierung der Seiten ab hier: 1, 2, 3, 4...
\chapter{Introduction}

Today radio resources become expensive due to the huge number of used systems. Smart ways to radio resource allocation is getting more and more important in wireless communication. The spectrum sharing ways should be done very carefully, to hold interference for legal systems. Additional approaches to increase frequency efficiency raise the complexity of the wireless systems \cite{Book32}. 
The fourth generation was the OFDM system. The OFDM system doesn't provide the necessary inter-system interference level for the nearby frequency bands \cite{Book34}\cite{Book35}. The OFDM system based on the orthogonality between different sub-carriers. The OFDM use additional time domain filter to archive orthogonality between sub-carriers.  In the next generation proposed more efficient way to use frequency resources. A new technique  is called generalized frequency division modulation. It based on the non-orthogonal sub-carrier location \cite{Book34} \cite{Book33}. The transmitter put the sub-carrier frequencies with overlap. The system provide such behaviour with root-raised cosine filters instead of the $sinc$ filters. Selected approach allow to decrease out-of-band radiation, but introduce self-interference in the system. The self-interference decrease performance in the system.  The new filter design add flexibility to the system because has the adjusting coefficient. The system adjust self-interference do find the trade-off between performance and frequency efficiency. Another requirement is spectrum sensing approach for the transmitter and receiver side. For that case receiver should estimate both the transmitted symbols and the used sub-carriers. This task has uncertainty in case if all symbols is unknown. To estimate both of the variables the receiver should know additional information. As example receiver may known the first column of the transmitted symbol matrix. All explained approaches make receiver and transmitter more complex. The receiver in the new system must solve much more complicated task to provide the same BER. The algorithms to symbol estimation also become more complicated.   There is a way to simplify the definition of the modulation matrix. We can use the tensor algebra and explain the transmission with a more natural and easy way. 
\section{Tensor based processing}\label{sec:TBP}
A tensor is a multidimensional array \eqref{i_1}\cite{Book6}. In other words, the Nth-order tensor is an matrix which has elements over N different dimensions. The first and second order tensors correspond to the arrays and matrices from linear algebra\cite{Book12}. The tensors which has number of dimension higher or equal to three are called higher order tensors. 
The order of the tensor is the number of dimensions which tensor has. In this thesis scalars are denoted by the lowercase letters e.g. a. The arrays are denoted by boldface lowercase letters e.g. $\mathbf{a}$. The matrices is written as the boldface capital letters e.g. $\mathbf{A}$. The higher order tensors are denoted by capital calligraphic letters e.g. $\mathcal{A}$.
Tensors explain data with high number of variables in a more natural way. Tensor based explanation of the data allow to consider all dependencies between dimensions of the data. 
The N-th order unfolding of the tensor is the stacked version of the tensor with respect to the N-th dimension \eqref{i_2}\eqref{i_3}\eqref{i_4}\eqref{i_5}\cite{Book16}. In the N-th order unfolding the elements of the each columns is fixed for the all dimensions except N. The N-th order unfolding has the number of row equal to the N dimension size, the number of the column equal to the product between other dimensions. The order of the columns are defines as the notation of the unfolding. In this paper is considered MATLAB ordering\cite{Book19}. The MATLAB notation change dimensions in the increasing order.
The N-th order tensor is rank-one if it can be presented as the outer products of N vectors.
The tensor decomposition is proposed by the Hitchcock in the 1927. He expressed the higher order tensor as the sum of the rank-one tensors \eqref{i_6}\eqref{i_7}.
The  tensor can be decomposed as the number of rank-one tensors. Decomposition called "CP" or CANDECOMP/PARAFAC decomposition\cite{Book6}. The "CP" decomposition factorizes the tensor into a sum of rank-one tensors components. The sum of rank-one tensor can be presented in another form \eqref{i_8} . The arrays corresponding to the same dimension can be stacked in the one matrix for each dimension respectively. This notation is called the "factor matrices". The factor matrices present the tensor unfolding in simplified form \eqref{i_10}\eqref{i_11}\eqref{i_12}.
There are many decomposition algorithms. One of the most popular is the Alternating least squares algorithm. 

\subsection{PARATUCK2 model}\label{part:PTM}
The PARATUCK2 model is the combination between the PARAFAC \cite{Book12} and TUCKER2 \cite{Book6} tensor models. The PARATUCK2 \cite{Book6} combines decomposition between five different matrices into one tensor. The
PARATUCK2 model in general is explained as the three dimensional
tensor and is defined in scalar form as given in \eqref{p_1} \cite{Book26}
where $x_{i_1,i_2,t}$ is $(i_1,i_2,t)$ -th entry of the resulting tensor $\mathcal{X} \in \compl^{I_1 \times I_2 \times T }$.The $\mathbf{S}$ matrix is the core matrix of the PARATUCK2 model. The matrices $\mathbf{A}$ and $\mathbf{B}$ show connection
between the core matrix $\mathbf{S}$ to the first and the second dimensions. The $\mathbf{C^{[a]}}$ and $\mathbf{C^{[b]}}$ are the scaling matrices defining the $\mathbf{A}$ and $\mathbf{B}$ relations to the core matrix $\mathbf{S}$.  Every element of the core matrix links the $\mathbf{A}$ and $\mathbf{B}$ matrices as the scaling coefficient\cite{Book12}.

The PARATUCK2 model possible to write in the two forms: slice-by-slice multiplication form and vectorized form.
In the  slice-by-slice product form the equation for the tensor is written for each slice separately \eqref{p_2}. Explained form allow to write the model as matrix multiplication with change only in two matrices for different slices.  The $\mathbf{C^{[a]}}$ and $\mathbf{C^{[b]}}$ is written as the tensor with elements on the diagonal in the first and second dimension. The linear combination between tensors become the scaling coefficients for columns in the each $i_3$ point. In that case the PARATUCK2 model will be written as the slice-wise multiplication \eqref{p_2} over the matrices $\mathbf{A} $ and $\mathbf{B} $ from the corresponding side with diagonalized tensors $\mathbf{C^{[a]}}$ and $\mathbf{C^{[b]}}$ for the corresponding slice\cite{Book12}.
Vectorized form of tensor $\mathcal{X}$ is another form to write the PARATUCK2 model\eqref{p_3} . We can write three equations with respect to the vectorized versions of the $\mathbf{A}$ \eqref{p_4},$\mathbf{S}$\eqref{p_3},$\mathbf{B}$\eqref{p_5}\cite{Book26}.Structure of the tensor in the defined equation is lost,  because we have only vectorized version of the tensor. The slice-wise model is replaced via matrix operations and allow to operate with matrices instead of the tensors of different matrix operation for different slices\cite{Book26}.
 The PARATUCK2 model in vectorised form possible to simplify with the Hadamard and the Khatri-Rao product notation in case if $\mathbf{A}$ and $\mathbf{B}$ have only one dimension \eqref{i_2}. The definition is provided via vectorization over the resulting tensor.The PARATUCK2 model become array in the third dimension.    The PARATUCK2 with one dimensioned $\mathbf{a}$ and $\mathbf{b}$  model is explained in the matrix notation with two different equation  \eqref{p_6}\cite{Book26}. The difference between equation is order of the multiplication between matrices. The Khatri-Rao and the Hadamard product is sensible to the multiplication order and change of the multiplication order, change overall model construction. The transpose of the $\mathbf{X}$ also change the matrices multiplication order.

The explained PARATUCK2 model will be used in the further work to define the GFDM system and allow to calculate solutions in the more natural form.
\section{Generalized frequency division modulation systems}\label{sec:GFDM}
The generalized frequency division modulation (GFDM) is the new approach to increase the spectral efficiency of the data transmission systems\cite{Book25} \cite{Book22}. The GFDM is extended version of the orthogonal frequency division modulation system with the additional overlapping in the frequency domain. The main principle of the GFDM system is transmission of the symbols during one transmission block on the number of sub-carriers and in the number of the time slots. Each symbol corresponding to the certain time slot is modulated with time domain filter. Time domain filter provides certain envelope in the time domain. The GFDM system use special form of the time domain filters which has the overlap with the symbols on the nearest sub-carriers. This approach provide more dense symbol arrangement in the frequency domain\cite{Book24}.
\subsection{Time filter construction principles}\label{part:TFCP}
 In the GFDM system are used two main time domain filters: raised cosine and root-raised cosine filter instead of the $sinc$ pulse shaping filter in the OFDM system.  Both of the filters have the variable coefficient $\alpha$ which allow to adjust the overlapping in the frequency domain from $0$ to $1$ where when $\alpha=0$ the overlap is minimal and  when $\alpha=1$ the overlap is maximal\cite{Book30}\cite{Book21}. That variable increase significantly inter-symbol interference and inter-channel interference and allow to put symbols more dense than in the OFDM system. The additional price for efficient frequency usage is receiver complexity. The receiver must separate all time symbols inside the one transmission block, even if the symbols is overlapped. 
 
 The raised cosine filter envelope is provided in the \eqref{g_1} below from the \cite{Book5} paper. In the \eqref{g_1} $T$ is the symbol duration in the time samples, and $\alpha$ is $roll-off$ factor, which is described above\cite{Book33}. The envelope of the raised cosine filter with different roll of factors is presented in the \ref{fg_1}. The following frequency envelope is presented in the \ref{fg_2}.

The comparison between the  $sinc$, raised cosine and root-raised cosine filter is shown in the fig. \ref{fg_5}. As we can see the overlap in the frequency domain spread energy from the symbol over the higher frequency range than for the $sinc$ filter \ref{fg_6}.

%\section{Iterative soft threshold optimization algorithm}\label{sec:ISTA}
%One of the most frequently solved tasks is deconvolution problem. The problem in the original form is written in the following equation \eqref{ista_1}. There are many ways to solve this problem in the second norm sense, with regularization and others additional algorithms. The problem in the deconvolution problem is very high conditional number of the matrix $\mathbf{H}$ which increase significantly the noise influence in the solution\cite{Book28}.
%\begin{align}
%R(x)=\mid\mid\mathbf{y}-\mathbf{Ax}\mid\mid^2_2 
%\label{ista_1}
%\end{align}
%\begin{align}
%R(x)=(\mathbf{y}-\mathbf{Ax})^H(\mathbf{y}-\mathbf{Ax})
%\label{ista_2}
%\end{align}
% In the most usual case algorithm should find the matrix inverse or pseudo inverse to solve the problem. In the paper \cite{Book27} was proposed another solution. We can find the iterative solution which will certainly decrease residual between the original and the reconstructed data from iteration to iteration. The main approach used in the algorithm called majoriziation-minimization \cite{Book27}. The approach replace the difficult problem by a number of easier minimization problems. Each minimization problem provide the answer for the solution, each next answer will be the closer to the right solution. We add additional non negative value to the original objective function to make problem easier. There are certain requirements for the majorization function $J_{add}(\mathbf{x})$. The $J_{add}(\mathbf{x})$ function must satisfy equation $J_{add}(\mathbf{x})\geq R(\mathbf{x})$ for all $x$ and moreover function must be equal in the certain point $J_{add}(\mathbf{x_k})= R(\mathbf{x_k})$. The function $J_{add}(\mathbf{x})$ can be different in the each certain point $\mathbf{x_k}$ and must be minimized easily. The explained below approach called "The Lanweber iteration". In the original form it was written for the real values function but we extend the algorithm for the complex valued functions. Assume the $J_{add}$ defined in the \eqref{ista_6}. In that case $ R_1(\mathbf{x})$ will be written as \eqref{ista_7} where $\alpha$ is some constant with is equal or higher to the maximal eigenvalue of the $\mathbf{A}^H\mathbf{A}$.  Rewrite the equation in the simpler form by opening the brackets \eqref{ista_10}. 
% \begin{align}
% R(\mathbf{x_{k+1}})<R(\mathbf{x_k})
% \label{ista_3}
% \end{align}
%\begin{align}
% R_1(\mathbf{x})=R(\mathbf{x})+J_{add}(\mathbf{x}) \label{ista_4}
%\end{align}
%\begin{align}
% R_1(\mathbf{x})=\mid\mid\mathbf{y}-\mathbf{Ax}\mid\mid^2_2+J_{add}(\mathbf{x})\label{ista_5}
%\end{align}
%\begin{align}
%J_{add}(\mathbf{x})=(\mathbf{x-x_k})^T(\alpha \mathbf{I}-\mathbf{A}^T\mathbf{A})(\mathbf{x-x_k})\label{ista_6}
%\end{align}
%\begin{align}
% R_1(\mathbf{x})=\mid\mid\mathbf{y}-\mathbf{Ax}\mid\mid^2_2+(\mathbf{x-x_k})^H(\alpha \mathbf{I}-\mathbf{A}^H\mathbf{A})(\mathbf{x-x_k}) \label{ista_7}
%\end{align}
%\begin{align}
% R_1(\mathbf{x})=(\mathbf{y}-\mathbf{Ax})^H(\mathbf{y}-\mathbf{Ax})+(\mathbf{x}-\mathbf{x}_k)^H(\alpha \mathbf{I}-\mathbf{A}^H\mathbf{A})(\mathbf{x}-\mathbf{x}_k) \label{ista_8}
%\end{align}
%\begin{align}
%\alpha \geq max(eig(\mathbf{A}^H\mathbf{A})) \label{ista_9}
%\end{align}
%\begin{align}
%R_1(\mathbf{x})=\mathbf{y}^H\mathbf{y} -\mathbf{y}^H\mathbf{A}\mathbf{x} -\mathbf{x}^H \mathbf{A}^H\mathbf{y} +\mathbf{x}^H\mathbf{A}^H\mathbf{A}\mathbf{x} \label{ista_10}
%\end{align}
%\begin{align*}
% +\mathbf{x}_k^H(\alpha \mathbf{I}-\mathbf{A}^H\mathbf{A})\mathbf{x}_k +\mathbf{x}^H(\alpha \mathbf{I}-\mathbf{A}^H\mathbf{A})\mathbf{x} 
%\end{align*}
%\begin{align*}
%-\mathbf{x}_k^H(\alpha \mathbf{I}-\mathbf{A}^H\mathbf{A})\mathbf{x} -\mathbf{x}^H(\alpha \mathbf{I}-\mathbf{A}^H\mathbf{A})\mathbf{x}_k 
%\end{align*}
%Equate the partial derivative to the zero \eqref{ista_12}. Because of the non holomorphic objective function, we take derivative with respect to the $\mathbf{x}^*$. During the transformations we can define \eqref{ista_14} the $\mathbf{x}$ from the known variables and explicitly write the equation to change $\mathbf{x}_k$ from iteration to iteration \eqref{ista_15}.
%\begin{align}
%\frac{\delta R_1(x)}{\delta x^*}=-\mathbf{A}^H\mathbf{y} +\mathbf{A}^H\mathbf{A}\mathbf{x} +(\alpha \mathbf{I}-\mathbf{A}^H\mathbf{A})\mathbf{x} -(\alpha \mathbf{I}-\mathbf{A}^H\mathbf{A})\mathbf{x}_k \label{ista_11}
%\end{align}
%\begin{align}
%\frac{\delta R_1(x)}{\delta x^*}=-\mathbf{A}^H\mathbf{y} +\alpha\mathbf{x} -\alpha\mathbf{x}_k+\mathbf{A}^H\mathbf{A}\mathbf{x}_k=0 \label{ista_12}
%\end{align}
%\begin{align}
%\mathbf{A}^H(\mathbf{y} -\mathbf{A}\mathbf{x}_k) +\alpha\mathbf{x}_k =\alpha\mathbf{x} \label{ista_13}
%\end{align}
%\begin{align}
%\mathbf{x}=\frac{\mathbf{A}^H}{\alpha}(\mathbf{y} -\mathbf{A}\mathbf{x}_k) +\mathbf{x}_k \label{ista_14}
%\end{align}
%\begin{align}
%\mathbf{x}_{k+1}=\frac{\mathbf{A}^H}{\alpha}(\mathbf{y} -\mathbf{A}\mathbf{x}_k) +\mathbf{x}_k \label{ista_15}
%\end{align}
%The defined equation \eqref{ista_15} allow to us write the optimization algorithm with high convergence speed and low computational complexity. Also there is one modification of the ISTA called FISTA which allow increase convergence speed for the one order.
%The FISTA algorithm comparing to the ISTA is similar to the CG algorithm in comparison with the SD\cite{Book28}. The FISTA works in the few steps:
%\begin{itemize}
%\item Set the start point $\mathbf{x}_0$ and set additional variable $t_0=1,t_1=1$
%\item Calculate $t_1$ via \eqref{ista_16}
%\item Take $y_1=x_0$
%\item Solve equation \eqref{ista_17} for the first iteration.
%\item Put $t_0=t_1$
%\item $y_{k+1}=x_k+(\frac{t_0-1}{t_1})(x_k-x_{k-1})$
%\item Check is the step size less than tolerance if no, go to step 2
%\end{itemize}
%\begin{align}
%t_1=\frac{1+\sqrt{1+4t_1}}{2} \label{ista_16}
%\end{align}
%\begin{align}
%\mathbf{x}_{k+1}=\frac{\mathbf{A}^H}{\alpha}(\mathbf{Y} -\mathbf{A}\mathbf{y}_k) +\mathbf{y}_k \label{ista_17}
%\end{align}
%The explained FISTA algorithm is used in the presented work in the system to solve ill-conditioned tasks for the ALS algorithm.
\chapter{Generalized frequency division multiplexing single input single output}\label{sec:GFDMSISO}
\section{System model}\label{part:SMSISO}
In this subsection will be explained the system model for the GFDM . We consider the system where transmitted signal is sent through the channel. The channel in physical meaning is the number of multipath propagation components from transmitter to the receiver or the convolution of the transmitted signal with some channel profile.  We  explain the received signal as matrix product between convolution matrix and the transmitted data array \eqref{siso_1}, where $\mathbf{y}$ is received array, $\mathbf{x}$ is transmitted array and $\mathbf{H}$ is channel convolution matrix. Also the additive white Gaussian noise is assumed. The  different channel models are considered in the thesis\eqref{siso_2}\eqref{siso_3}\cite{Book21}.
The transmitted data is the third order PARATUCK2 tensor in the vectorized or unfolding form. The symbol matrix $\mathbf{S}$ introduced as the core matrix in the PARATUCK2 model\cite{Book6}. The symbols are written in the core matrix without any correction.
 Each column in the $\mathbf{C^{[a]}}$ matrix corresponds to the certain sub-carrier frequency for the stream modulation\eqref{siso_4}\eqref{siso_5}.The $\mathbf{C^{[a]}}$ matrix is constructed as the sub-carrier frequency  time samples in the each column, where column $i$ will correspond to the $i$ sub-carrier in the GFDM system. The overall equation to the $\mathbf{C^{[a]}}$ construction is explained in the \eqref{siso_1}. 
 Each column in the $\mathbf{C^{[b]}}$ define spread of the certain symbol over the all transmission block \eqref{siso_6}\cite{Book12}.  Elements in the columns are the time filter functions for the symbols. Each column is the shifted version of the other. The shift between near columns is equal to the $ T/T_s$ \eqref{siso_7}. The time filters functions are the root-raised cosine filter with adjustable $\alpha$ coefficient.

 The  $\mathbf{B}$ matrix become $\mathbf{b}$ array
and in the physical meaning corresponds to the time variation inside the one transmission block for different symbols or as selection coefficients for the different time slots\cite{Book34}.  We assume that system is linear time invariant and therefore $\mathbf{b}$ become array with all-ones structure \eqref{siso_8}.

The  $\mathbf{A}$ matrix become $\mathbf{a}$ array and in the physical meaning corresponds to the transmission coefficient for each sub-carrier frequency stream or  selection array for sub-carriers inside of the one transmission block \eqref{siso_10}. We will consider the system where the $\mathbf{a}$ will has two typical structures. First is the array with all-ones in case if there is no overlapping in the frequency domain systems and known to the receiver. In the second case the values in the $\mathbf{a}$ is unknown for the receiver and have random variables from the range $[0,1]$

The transmitted data is defined in equation \eqref{siso_12} where each matrix is explained above. Regarding to the chapter X the PARATUCK2 model can be defined in the different equations \eqref{siso_13} \eqref{siso_14} with vectorization operation for the resulting tensor. The vectorized model is sensible to transpose of the resulting matrix and to the order change.

\section{Additive white Gaussian noise channel}\label{part:AWGN}
In this section we assume the GFDM transmitter-receiver pair in the SISO case with the additive white Gaussian noise at the receiver and memoryless flat-frequency channel\cite{Book30}. Therefore system model at the receiver become equation \eqref{siso_15}.

\subsection{Zero forcing solution}
The receiver can find the optimal transmitted values via optimization task solution\cite{Book11}.  The demodulation problem can be rewritten in the residual form. We can rewrite $\mathbf{y}$ as the product of the known modulation matrix with the vectorized symbols matrix \eqref{siso_ls_1}. We introduce the matrix $\mathbf{\Omega}_1$ in \eqref{siso_ls_2} and define $vec(\mathbf{y)}$ as the matrix product between the $\mathbf{\Omega}_1$ and estimated $vec(\mathbf{\widehat{S}})$ \eqref{siso_ls_3}. We define residual equation which is difference between received data and the estimated at the receiver symbols \eqref{siso_ls_4}. The  $\mathbf{\Omega}_1$ is given from the paper \cite{Book26}. If the $\mathbf{a}$ and $\mathbf{b}$ have all-ones structure we can neglect them and take as the Khatri-Rao product between $\mathbf{C^{[a]}}$ and $\mathbf{C^{[b]}}$ \eqref{siso_ls_4}
The receiver must solve the system of linear equations. The problem can be solved in the second norm sense \eqref{siso_ls_5} and also known as the Zero-Forced receiver \cite{Book24}.The second norm function is convex, with one global minimum point. To find minimum point the receiver find partial derivative of the objective function \eqref{siso_ls_6}, equate it to the zero and solve the system of linear equations. The estimated $\mathbf{S}$ matrix is complex valued, which makes objective function non-holomorphic. To solve optimization problem we have used the Wirtinger calculus and find the partial derivative with respect to the $vec(\mathbf{\widehat{S}}^*)$ and equate it to the zero. Next calculations steps come to the classical  least squares solution of the objective function \eqref{siso_ls_10}\cite{Book24}. 
\subsection{Matched filter solution}
Another solution for the problem is so called Matched filter receiver. The receiver  multiply received data with the hermitian transpose of the modulation matrix\cite{Book25}. We can look at the singular value decomposition of the $\mathbf{\Omega}_1$ \eqref{siso_mf_1} and write in the same form the inverse of the matrix due to the unitary structure of the matrix $\mathbf{U}$ \eqref{siso_mf_2}. The inverse of the $\mathbf{\Sigma}$ is inverse of the each element of the matrix due to the diagonal structure\cite{Book25}. With increasing of the $\alpha$ coefficient in the time filter the condition number of the $\mathbf{\Omega}_1$\eqref{siso_mf_4} decrease and inverse of the matrix will increase noise components with decreasing minimal singular value of the $\mathbf{\Omega}_1$\eqref{siso_mf_4}. We can multiply the received data with hermitian of the $\mathbf{\Omega}_1$\eqref{siso_mf_5}\cite{Book24}. In that case the structure of the singular components inside of the  calculated matrix become the same and noise components will have the same power\eqref{siso_mf_6}.

\section{Selection coefficients search}\label{part:SCSSISO}
Explained above the coefficients from the $\mathbf{a}$ are used as the selection coefficients for the sub-carriers. In the GFDM system is allowed to disable some sub-carriers to decrease interference with other systems which work in the same frequency range\cite{Book25}.The transmitted may use the spectrum sensing approach and doesn't use the number of subcarriers inside of the transmission block. Receiver will not know which sub-carriers are used and must estimate information from the received data\eqref{siso_s_1}. The information about used sub-carriers may transmitted as additional to find coefficients at the receiver side.
We can use the PARATUCK2 model to estimate information about used subcarriers at the receiver. The received signal is defined as \eqref{siso_s_1}, where the selection coefficients become in the $\mathbf{a}$

We formulate equation to vectorize the $\mathbf{a}$ from the right hand side as \eqref{siso_s_2} and write the problem statement to find the coefficients as the residual between received signal and estimated channel \eqref{siso_s_3}.
\subsection{Semi-blind receiver}
The receiver estimate the transmission coefficients and the symbols in case if will know at least one symbol at each sub-carrier frequency. The receiver will estimate in the certain time slot the transmission coefficient for the overall block of the data and will find all other symbols according to the LS solution or optimization algorithm. We should write the objective function and minimize it to implement the semi-blind receiver. The unknown variables are $\mathbf{\widehat{S}}$ and the $\mathbf{a}$. The known symbols is added in the objective function as constraint or via decreasing number of the variables.
 Second way decreases scale of the objective function, but problem formulation  become non-trivial. We add known symbol information as the additional constraint in the objective function  \eqref{siso_sm_1}. There are two ways to solve the task, joint and  approximated solutions. The joint solution consider  all dependences between $r_1,r_2,r_3$. The second approach calculate solution for three residual function separately and use non-linear optimization process to stack three equations in one. We implemented both of them.
\subsection{Approximated semi-blind receiver}
 Problem statement is defined in \eqref{siso_sm_2}. We minimize objective function with assumption: $\mathbf{r}_1$ and $\mathbf{r}_2$ does not depend from each other.  We don't used the Lagrangian multipliers to analytical optimization and find minimum point in the optimization process. The residual function $r_1$ corresponds to the objective function with vectorized symbol matrix $\mathbf{\widehat{S}}$. The $\mathbf{\Omega}_1$ is defined previously and corresponds to the modulation matrix.
 The residual function $r_2$ corresponds to the objective function with vectorized array $\mathbf{a}$. The $\mathbf{\Omega}_2$ is defined previously and corresponds to the transmitted data with separated sub-carrier coefficients.
The residual function $r_3$ corresponds to the constraint for the known symbols, where $\mathbf{S}_{sel}$ become the selection matrix, where on the main diagonal is ones in points where $vec(\mathbf{\widehat{S}})$ is known. The $\mathbf{q}$ array is the array with vectorized known symbols in the $vec(\mathbf{\widehat{S}})$ and zeros in the unknown points.
We can stack \eqref{siso_sm_10}\eqref{siso_sm_9}\eqref{siso_sm_8} in one system of non-linear equations.There are many ways to sole the system of non-linear equations.
 There are two additional approaches which could be implemented in the Newton algorithm. One is the Powell-Wolf rule\cite{Book62}\cite{Book63} for the step-size variable adjust. The second is the Levenberg-Marquardt regularization\cite{Book66}\cite{Book68} algorithm. The main equation of the Newton algorithm for the non-linear equations is presented in the \eqref{siso_sm_11}, where $\mathbf{d}$ is the derivative array which is minimized to the zero and $J$ is the Jacobian matrix of the $\mathbf{d}$ array. The definition of the partial derivative array $\mathbf{d}$ array is  explained in the equation\eqref{siso_sm_13}. 
We must calculate partial derivative from the  $\mathbf{d}$ to find the Jacobian matrix \eqref{siso_sm_18}.Because the function\eqref{siso_sm_13} is holomorphic, we don't use the Wirtinger calculus in Jacobian calculation. We calculate partial derivative with respect to the $\delta \theta$ \eqref{siso_sm_19}. 
 The receiver solve the system of the linear equations at the each iteration and update the search vector for the next step\eqref{siso_sm_12}. The optimization algorithm decrease the derivative of residual to the zero, and converge due to the convexity of the second norm\cite{Book58}.

\subsection{Joint semi-blind receiver}
We can minimize the objective function \eqref{siso_jsm_1} without assumptions and approximations. General form makes implementation more complex. It should be noticed that $\mathbf{r}_1$ and the $\mathbf{r}_2$ is the same equations but explained from the different point of view: from the unknown symbols or unknown sub-carrier coefficients.  
We write the objective function. The function in non-holomorphic. We find the partial derivatives with Wirtinger calculus. We can use the \eqref{siso_jsm_3} property of the derivatives to find the derivatives separately for each function. To calculate the cross derivatives for the $\mathbf{r}_1^H\mathbf{r}_1$ with respect to the $\mathbf{a}^*$ we use the \eqref{siso_jsm_7} which allow to change the left hand part of equation, where $\mathbf{\widehat{a}}^*$ is clearly defined\cite{Book26}.
The optimization function become the system of non-linear equations. As explained previously, to solve the system of the non-linear equations we use the Newton method.The $\mathbf{d}$ is presented in the \eqref{siso_jsm_12}. The function is holomorphic\cite{Book61} and we can calculate Jacobian matrix  $\mathbf{J}$ without Wirtinger calculus \eqref{siso_jsm_11}. The $r_1$ and the $r_2$ is equal function in the different notation. We use this property when we calculate cross variable derivatives.  The resulting Jacobian matrix \eqref{siso_jsm_13} is different from the Jacobian matrix presented in the approximated solution \eqref{siso_jsm_13}

Joint solution converge faster in comparison with the approximated approach. Each iteration in the joint solution algorithm become computationally more expensive, but algorithm takes much less iterations.
\section{Frequency selective channel estimation}\label{part:FSCESISO}
\subsection{Approximated channel estimation}
In this section is considered model where channel has the multipath propagation components \eqref{siso_a_1}\eqref{siso_a_2}. We use assumption of maximal delay between first and the last multipath propagation components equal to the $T/T_s$ time samples. This assumption makes approach below only approximation of the channel orthogonalization .
We can implement the channel estimation approach without channel prefix in the start of the transmission block. We put the training symbol in the transmission block. 
Consider the matrix $\mathbf{H}$ in terms of parametric model. The matrix $\mathbf{H}$ has Toeplitz lower triangular structure\eqref{siso_a_3}\cite{Book5}. Each column in the $\mathbf{H}$ corresponds to the multipath components  for certain time sample. We assume that length of the channel is less than $T/T_s$. The matrix $\mathbf{H}$ has only $T/T_s$ non-zero elements in the each column. Maximal length of the cyclic prefix to estimate the channel equal to the $T/T_s$ samples in that case. 
The transmitter insert this prefix as predefined symbols in the first column in the matrix $\mathbf{S}$ \eqref{siso_a_5}. The first column in the matrix $\mathbf{S}$ corresponds to the first time slot for each sub-carrier stream\eqref{siso_a_6}. The knowledge of the symbols in the first column doesn't define the first $T/T_s$ time samples in the transmission block. The receiver doesn't know exactly the first $T/T_s$ samples of received data, even if the first $F$ symbols is known. Due to that fact the receiver will approximate the channel with the error \eqref{siso_a_13}\cite{Book23}. Define the matrix with training symbols\eqref{siso_a_5}.Define the matrix with unknown  symbols\eqref{siso_a_6}. The training signal from the transmitter will be written in the following form \eqref{}. The received signal will be written as following\eqref{}. The receiver can construct the following array from the known training symbols\eqref{} and define the convolution matrix from the known array. 

One approach to decrease channel estimation error is the adaptive $\alpha$ coefficient adjustment. The transmitter use two different $\alpha$ coefficients for the working time slots and for the test time slots. We use guard interval with length one symbol to avoid ICI and decrease interference for the certain symbols. Symbols with decreased ICI used as the cyclic prefix and increase estimation precision for the described above approach. 
The second approach to decrease uncertainty and increase precision is optimization algorithm. Algorithm estimate use additional double side interference cancellation. The joint symbol and multipath components estimation leads to the decrease of uncertainty with symbol estimation, and decrease probability of the error with the multipath component estimation. 
\subsection{Channel estimation}
The main equation of the received data \eqref{ce_2} can be rewritten with respect to the unknown channel model \eqref{ce_1}.

Where the $\mathbf{h}$ is the channel coefficients with known length $L+1$. The matrix $\mathbf{D}_{mod}$ is constructed via shifting and stacking of the transposed array $\mathbf{\Omega}_1vec(\mathbf{S})$\cite{Book53}.
As shown in the \eqref{ce_4} the receiver can estimate the channel is case if symbols are known\cite{Book47}. It should be noted that transmitter must know the length of the channel . The least squares solution is calculated as following in the previous parts with the Wirtinger calculus \eqref{ce_5}. The receiver calculate the pseudo-inverse of the matrix \eqref{ce_7}. The solution for the least squares problem is presented in the \eqref{ce_7}.

 In the practical case the transmission block is too large to them all with known symbols. The transmitter can fill number of the symbols with the known data and other as unknown to the receiver. The matrices $\mathbf{S}_{sel}$ and $\mathbf{S}_{sel unk}$ are constructed in the similar way. The matrices are the selection matrices with ones in the points where symbol transmitted with certain column of the $\mathbf{\Omega}_1$.


 The receiver must find both symbols and the channel values and perform the semi-blind receiver solution. To make the statement of the task separated from the known and unknown symbol part, which also decrease inter-symbol interference the receiver rewrite the task in the following notation \eqref{ce_17}.

 We write the received data as the sum of two array: known and unknown. The receiver also can divide the modulation matrix into the two parts: with columns corresponding to the known symbols and with other columns of the matrix $\mathbf{\Omega}_1$.
  With such equation the unknown symbols will estimated more precisely in comparison with simple ALS approach. The receiver construct the two residual equations and try to minimize them during the optimization task. The residual function based on the second norm to make the objective function convex. The optimization task can be solved with ALS and Newton algorithm. 
 
  The objective function is following\eqref{ce_als_1}.  We construct the partial derivative with respect to the unknown symbols and channel values. We used the Wirtinger calculus to find the derivative of the complex function. The derivative is written as the system of non-linear equations. We used ALS and Newton algorithms to solve the problem . 
 

The Newton algorithm solve the system of linear equations at each iteration for which given point \cite{Book62}. We must write the equation of the Jacobian matrix \eqref{ce_n_3} for the  partial derivatives to determine the solved at each iteration equation. To write the Jacobian matrix we use property of the equality of the different notations of the same function in the same point. The resulting Jacobian matrix is written as follows \eqref{ce_n_5}. 
The Newton algorithm may be regularized and written in the more stable form with the Powell-Wolf stepsize rule\cite{Book66} and Levenberg-Marquardt regularization approach \cite{Book65}. 
%It is also possible to use explained in the \cite{Book53}\cite{Book52} so called semi-blind receiver based regularization approached, but they request the estimation over number of the received blocks of the data.
\section{Simulation results}\label{part:SRSISO}
In this section we present simulation results for the GFDM system and algorithms which was reviewed in the chapter \ref{sec:GFDMSISO}.

Performance of the GFDM system in comparison with different $\alpha$ coefficients are obtained through simulation. The parameters of the system are tabulated in Table \eqref{tab:table1}. The GFDM system is simulated in the AWGN channel without coding. QPSK modulation scheme is used. The number of sub-carriers was $F=32$, samples for each symbol was $T/T_s=F$ The block size was $T_s=15$. The root-raised cosine filter was used with  list of roll-of-factors. In the roll-of factor test we measured the SER for the different $\alpha$ for the Zero-Forced receiver and matched filter. The resulting SER plot is shown for the zero forced receiver at fig. \eqref{fs_1} and  for the matched filter receiver at fig. \eqref{fs_2}.

Performance results of the GFDM selection coefficient estimation and semi-blind receiver are obtained through simulation. The parameters of the system are tabulated in Table \ref{tab:table2}. 

The GFDM system is simulated with AWGN channel. QPSK modulation scheme is used. The number of sub-carriers was $F=32$, samples for each symbol was $T/T_s=F$ The block size was $T_s=15$. The root-raised cosine filter was used with roll-of-factor $\alpha=0.5$. The selection coefficients was chosen as the random integer values in range from $0$ to $1$.  The results of the GFDM performance is presented in the two plots, at first \eqref{fs_3} is shown the SER in comparison with the case  if receiver know the original $\mathbf{a}$ with the Zero-Forced receiver and with case if receiver doesn't know the $\mathbf{a}$ . At second figure \eqref{fs_4} is shown the reconstruction error for the $\mathbf{a}$ .

There are additional plots for the joint and approximated algorithms comparison is shown at  fig.\eqref{fs_5} fig.\eqref{fs_6} fig.\eqref{fs_7}. The comparison results are obtained through simulation. The setup of the simulation is shown from the Table \ref{tab:table3}.The GFDM system is simulated with AWGN channel. QPSK modulation scheme is used. The parameters of the setup is the same as in the previous experiment. In the fig.\eqref{fs_5} is shown the typical reconstruction error convergence of the both of the algorithms per iterations. As the real minimum point is shown values of the residual between generated matrices and received data. In the fig.\eqref{fs_6} shown the time of the convergence of the algorithm with dependence for the SNR. In the fig.\eqref{fs_7} shown the typical convergence of the residual \eqref{sim:eq_1} of the both of the algorithms per iterations. As the real minimum point is shown values of the residual between generated matrices and received data.

The GFDM system semi-blind receiver is simulated in the frequency selective channel without coding. QPSK modulation scheme is used. The number of sub-carriers was $F=8$, samples for each symbol was $T/T_s=F$ The block size was $T_s=3$. The root-raised cosine filter was used with roll-of-factor $\alpha=1$. The length of the channel equal to $L+1=4$. The values of the channel is presented in the table \ref{tab:sisotable5}. The results of the GFDM performance are presented in the two plots, at first \eqref{fs_8} is shown the SER in comparison with the Zero-Forcing FFT receiver and with case for different number of unknown symbols in one transmission block. In the simulation we define as the unknown symbol number of the non-training symbols in the $vec(\mathbf{S})$ At second figure \eqref{fs_9} is shown the reconstruction error for the channel.

\section{Conclusion}\label{part:CSISO}
As conclusion for the first experiment we  can say that zero forced receiver show much better results in the symbol detection in comparison with matched filter. The difference become significant when the $roll-off$ factor of the $RRC$ filter increase. The ZF receiver has increased the SER slightly with $\alpha$ increasing which is shown at the fig.\ref{fs_1}. The MF receiver increase SER significantly more  with $\alpha$ increasing which is shown at fig.\ref{fs_2}. Explained results show, that ZF receiver should be preferred in the receiver. The ZF receiver achieve the better performance with the same self-interference ratio and allow to decrease the out-of-band radiation.
As conclusion for the experiment with frequency coefficient selection approach we can see that algorithm decrease SER for 5-6 dB in comparison with system which know the frequency coefficients fig.\ref{fs_3}. This significant decrease in performance come from the non-convenient coefficient selection way in algorithm, where coefficient was chosen as $1$ if the resulting $abs(a_1)$ value is higher that 0.5. There are possible better solutions which increase performance of the algorithm and approaches to use predefined information at the receiver. We can see that joint algorithm show slightly worse performance in symbol and frequency coefficients detection fig.\ref{fs_3}fig.\ref{fs_4}in comparison with approximated algorithms. Those behaviour is explained from the approximated algorithm convergence. We take as the start point transmission coefficients equal to the zero. The approximated algorithm turn on less sub-carriers and make false negative errors. The false negative error lead to the better SER performance due to less estimated number of the symbols. The possible performance increase for the joint solution algorithm is the different weights for the objective function parts. The part for frequency coefficients estimations should have higher weight.
The third experiment shows advantages of the joint algorithm. As we can see from the all three figures \ref{fs_5},\ref{fs_6},\ref{fs_7}. The residual decrease show that joint solution algorithm converge much faster than approximated. The typical number of iteration of the joint algorithm equal to $2$ and for the joint algorithm equal to the $200$. So significant difference neglect the difference between the computationally expensive iterations in the joint algorithm. The fig. \ref{fs_6} confirm that fact. The joint algorithm converge independently from the SNR and take 100 times less time to converge. The approximated algorithm slightly depend from the SNR in the received data.
The fig. \ref{fs_7} show that estimated values in the approximated algorithm decrease non-linearly. The shown fact mean that algorithm decrease at each iteration only one set of the variables, symbols or frequency coefficients. The overall conclusion is that approach is effective enough for the data estimation and there are possibility to increase performance with weights addition. Algorithm will show slightly better results.
The next experiment is estimated the frequency selective channel estimation approach. The experiment was done for the all possible number of known symbols which have certain problem statement. The result for the SER if shown in the fig. \ref{fs_8}. As we can see here, there is one minimum point in symbol estimation. The less unknown symbols in the block leads to the SNR increase, whether the higher number of unknown symbols leads to the SER increase. The point where number of unknown symbols equal to $8$ show the closest result to the original GFDM system with the same block size but known channel and ZF receiver. The channel estimation error has the same behaviour and show the minimum point in the same number of unknown symbols. Important fact is that reconstruction error for the maximal and minimal number of unknown symbols have the same behaviour and lay in the same line. Which mean that maximum knowledge and maximal uncertainty show the same results. This behaviour also show the convexity of the performance function for the number of unknown symbols as variable. In the difference block size should be different optimal value of the unknown symbols, but the average performance must be the worse in comparison with the original GFDM system due to the higher energy per bit with same performance. There is additional important performance gain in the explained algorithm, the algorithm decrease the self-interference in the system from known symbols. The algorithm also can be written for the overall interference decrease, but better to use this algorithm after the semi-blind receiver.

\chapter{Generalized frequency division multiplexing multiple input multiple output}\label{sec:GFDMMIMO}
\section{System model}\label{part:SM1}
In this chapter we assume the GFDM MIMO system  the flat-frequency channel\cite{Book25}\cite{Book29} in that case the system is explained in the 
\eqref{mim_1} equation with channel model a identity matrix without multipath propagation.
%\subsection{MIMO model}
The GFDM system in MIMO AWGN case can be explained with the PARATUCK2 third order model. The difference between the SISO and MIMO case is significant\cite{Book29}, but PARATUCK2 model allow to extend  matrices to define the MIMO case inside of the same mathematical structure. The received data should have certain dimensionality $\mathbf{Y}\in\compl^{M_r\times T}$
The symbol matrix $\mathbf{S}$ from the SISO system become three dimensional tensor $\mathcal{S}\in\compl^{F\times T_s\times M_t}$ in the MIMO case. Additional dimension corresponds to the symbols which is transmitted from the different transmit antennas. The PARATUCK2 does not allow to use tensors in the model. To overcome this constraint we transform the $\mathcal{X}$ into the stacked slices for each $m_t$ \eqref{mimo_1}.
The $C^{[b]}$ matrix stay the same as in the SISO model and does not change with increasing of the transmit antennas \eqref{mim_2}. In the columns of the matrix stays the same time mixing filter which doesn't change from the number of the transmit and receive antennas. From the right hand side the $\mathbf{S}$ does not change and have number of the columns $T_s$ which also satisfy the model.

The connection array $\mathbf{b}$ stay the same as in the SISO case \eqref{mimo_3}. The necessary assumption for this is linear time invariant transmission channel in terms of the one transmission block. We fill the $\mathbf{b}$ array with all-ones as in the SISO case.
The $\mathbf{C}^{[a]}$ matrix in the MIMO case is changed \eqref{mimo_4}. Define this matrix in the MIMO case as $\mathbf{C}^{[a]'}$. Regarding to the $\mathbf{S}$ matrix, dimensions of the $\mathbf{C}^{[a]'}$ should be $T \times F\cdot M_t$. In comparison with the SISO model, in the MIMO case $\mathbf{S}$ should be multiplied with the $\mathbf{C}^{[a]}$ matrices. The new $\mathbf{C}^{[a]'}$ become stacked repeated version of the $\mathbf{C}^{[a]}$ matrix \eqref{mimo_5}. Hadamard product between $\mathbf{S}$ and the $\mathbf{C}^{[a]'}$ become simple explanation of the sub-carrier modulation of the each sub-frequency at the each transmitter antenna. Sub-carriers in the each transmit antennas is equal, so $\mathbf{C}^{[a]}$ matrices for each $m_t$ is also equal.
The $\mathbf{a}$ array become the matrix and connect each $m_r$ antenna with the each $m_t$ transmit antenna and include every sub-carrier coefficient\eqref{mimo_6}. Define sub-carrier coefficients as the $\mathcal{A}$ tensor with coefficients for each $m_r,m_t,f$. In the PARATUCK2 model $\mathbf{A}$ is the matrix, we must define the $\mathcal{A}$ as the matrix which will sum for every transmit antenna sub-carrier frequency with corresponding sub-carrier coefficient\eqref{mimo_7}. Regarding to the $\mathbf{S}$ matrix, symbols is located sequentially for each sub-carrier at each transmit antenna. The $\mathbf{A}$ dimension is $M_r \times F\cdot M_t$. The first unfolding of the $\mathcal{A}$ have corresponding sub-carrier coefficients following the $\mathbf{S}$ symbols\eqref{mimo_8}. This notation multiply each sub-carrier coefficient $\mathbf{A}$ to each symbol from the $\mathbf{S}$.
The overall model with explained matrices will be defined in the \eqref{mimo_9}. This model define the received data. All operation which was explained in the GFDM SISO part is also applicable in the GFDM MIMO case. It should be noted that the rank of the $\mathbf{A}$ is $min(M_r,F\cdot M_t)$ in the real case number of the receive antennas is much less than sub-carriers multiplied with transmit antennas. The rank of the matrix $\mathbf{A}$ show how many transmit coefficient possible to know in the system with known transmitted symbols.

\section{Sub-carrier coefficient update approach}\label{part:SCUA}
The spectrum sensing approach is applicable in the MIMO system too. One of the important properties in the MIMO system is diversity, which frequently achieved via spatial location of the transmitter and receiver antennas\cite{Book29}. It should be noted that for different transmitter-receiver antenna pair one sub-carrier may be occupied with another system, and not. Receiver must know the $\mathbf{A}$ to know how data is transmitted. Sub-carrier coefficients update is complex task due to the time mix and multi-frequency system. The number of the unknown variables in the system for MIMO case equal to the $M_r\cdot M_t\cdot F$. Number of the equations to solve equal to the $T\cdot M_r$ for each transmission block. The transmission coefficients can be found with specially structured known symbol matrix. 
\subsection{Approximated approach}
The one way to find the matrix $\mathbf{A}$ with sub-carrier coefficients is semi-blind receiver. The number of time slots must equal to the number of the transmit antennas. The receiver may use approach from the previous chapter \ref{part:SCSSISO}. Transmitter should construct special structure of the matrix $\mathbf{S}$. For each time slot the receiver can find the $M_r\cdot F$ transmission coefficients. The transmitter must fill at least $M_t$ time slots with known symbols to find the all $\mathbf{A}$ matrix. The search algorithm for the  matrix $\mathbf{A}$ is next:
\begin{itemize}
\item Transmitter put inside the transmission block specially structured symbols. Transmitter put known to the receiver symbol row to the one transmit antenna at the each time slot. In other words the symbols are sent from each sub-carrier frequency in the certain antenna.  Transmitter repeat the same row, but from the another transmit antenna at the next time slot. The transmitter repeat process $M_t$ times for each antenna. The receiver must know the   transmission order. 
\item  The receiver has transmission block with $M_t$ time slots length. The receiver uses semi-blind approach from the SISO case and assume received signal is sent from the SISO system. Semi-blind receiver approach estimates the sub-carrier coefficients for $m_t$-th transmit antenna.
\item The receiver repeats process for another time slot, and estimates the sub-carrier coefficients for another transmit antenna. The receiver repeats process $M_t$ times until estimates the all $\mathbf{A}$ matrix elements.
\end{itemize}
\subsection{MIMO model based approach}
The matrix $\mathbf{A}$ will be found with the equation \eqref{mimo_msm_4} separation  in the PARATUCK2 model. The model is written as the vectorized matrix $\mathbf{A}$ multiplication with the intermediate matrix $\mathbf{\Delta}$, as given in \eqref{mimo_msm_5}. 
 The product of this two matrices explained in the vectorized form for both of the second matrices \eqref{mimo_msm_4}.The matrix $\mathbf{\Delta}$ constructed with property \eqref{mimo_msm_4} and defined in \eqref{mimo_msm_5}.  We used property \eqref{mimo_msm_4} to find the sub-carrier coefficient matrix $\mathbf{A}$ from the received data. The overall equation is defined as \eqref{mimo_msm_6}.
 We constructed the objective function to minimize the second norm of the vectorized residual explained via $vec(\mathbf{A})$ array \eqref{mimo_msm_6}. The function is non-holomorphic\cite{Book58} due to the complex values in the  $vec(\mathbf{A})$. We have used Wirtinger calculus and have found the first partial derivative. We have equated first partial derivative to the zero\cite{Book47}. The first partial derivative has global minimum point, because the second norm  function is convex. 

 The optimal solution in the second norm sense is written as here \eqref{mimo_msm_15}. We can simplify defined equation \eqref{mimo_msm_15} with  linear algebra properties. The simplified equation is defined as \eqref{mimo_msm_17}

We applied the Kronecker product properties and have analysed the rank of the inverse matrix $((\mathbf{\Delta}^*\mathbf{\Delta}^T)^{-1} \otimes \mathbf{I}^{-1})(\mathbf{\Delta}^*\otimes \mathbf{I})$ to find number of the coefficients which is possible to estimate in the one transmission block \eqref{mimo_msm_17}. The rank of the Kronecker product equals to the product between the multiplied matrix ranks \eqref{mimo_msm_rank_1}. The identity matrix has the maximal rank, so the rank of the matrix is defined by the second part of the equation \eqref{mimo_msm_rank_3}. Consider the matrix \eqref{mimo_msm_5}, The size of the core matrix $\mathbf{S}$ is equal to the $M_t\cdot F\times T_s$ , rank of the matrix $(\mathbf{S}\cdot (\mathbf{C}^{[b]}\diamond \mathbf{B})^T)$ is less of equal to the $T_s$. There is inequality which define the rank of the Hadamard product \eqref{mimo_msm_rank_5} as written in the \cite{Book48}\cite{Book56}\cite{Book57}. The matrix $\mathbf{C}^{[a]'T}$ has the rank equal to the $min(F,T)$ because it based on the stacked repetition of the $\mathbf{C}^{[a]T}$ matrix. The $\mathbf{C}^{[a]T}$ matrix has the full rank $min(F,T)$. The matrix \eqref{mimo_msm_rank_6} has the rank equal to the matrix $\mathbf{\Delta}$ rank. We have defined the system of the linear equations rank with equation \eqref{mimo_msm_rank_7}.

The transmitter can use the specially structured symbol matrix as explained in the semi-blind receiver for SISO case part. We will find the all matrix $\mathbf{A}$ via least squares solution.We consider all dependencies with the MIMO based approach. We don't make assumption about SISO case and hold precision of the estimation.
\subsection{Decomposition based approach}
The receiver can estimate sub-carrier selection coefficients with less number of unknowns. The tensor $\mathcal{A}$ from the PARATUCK2 model is explained as unfolding of tensor decomposition and decrease number of estimated variables.

where $\mathbf{A}_A \in\compl^{M_r\times r} \mathbf{A}_B \in\compl^{F\times r} \mathbf{A}_C \in\compl^{M_t\times r}$ and therefore receiver decrease number of variables to solve the task\eqref{mimo_dsm_5}. Relation from $r$ to possible decrease of necessary equations grows polynomially. Rank of the decomposition decrease number of the variables in case if $r$ from the \eqref{mimo_dsm_7} will be higher than rank of the tensor decomposition.

We have written equation in the vectorized residual form and have classical system of linear equations. 

Possible to separate $vec()$ operation for the $\mathbf{A}_A$ and another part of the multiplication exploiting $vec$ operation properties \eqref{mimo_dsm_11} which we can see from the equation \eqref{mimo_dsm_8} .  The receiver estimate the $\mathbf{A}_A$ matrix in the least squares sense from equation \eqref{mimo_dsm_15}.

 The receiver separate Khatri-Rao product components from the equation \eqref{mimo_dsm_c_2}  using another property of the $vec$ operation \eqref{mimo_dsm_c_1}. The receiver can find Khatri-Rao product and must separate $\mathbf{A}_B$ and the $\mathbf{A}_C$ matrix from the product.
Possible to separate $\mathbf{A}_B$ and the $\mathbf{A}_C$ matrices from the Khatri-Rao product, using the decomposition of the Khatri-Rao product\eqref{mimo_dsm_c_11}. The receiver doesn't estimate over the Khatri-Rao product, but estimates the vectorized product version \eqref{mimo_dsm_c_5}. The receiver separate Khatri-Rao product components with the Hadamard product, but doesn't reveal necessary matrices in the matrix algebra. We introduce two other equations \eqref{mimo_dsm_c_12} \eqref{mimo_dsm_c_14} to reveal $\mathbf{A}_B$ and the $\mathbf{A}_C$ in the clear form. 

Matrices inside of the Khatri-Rao product is separable with vectorization operation following the equations \eqref{mimo_dsm_c_11} \eqref{mimo_dsm_c_12}. We can estimate each of the matrix and decrease number of the variables in comparison with matrix based approach.  We can separate the \eqref{mimo_dsm_c_11} with additional  presented operations.
We have constructed the objective function to estimate minimum point of the residual between estimated and the received data\eqref{mimo_dsm_opt_1}. We have rewritten the objective function with respect to the each of the variable set\eqref{mimo_dsm_opt_3}. The objective function is non-holomorphic \cite{Book52} and we have used the Wirtinger calculus to estimate the partial derivatives\eqref{mimo_dsm_opt_9} from the each of the variables set \eqref{mimo_dsm_opt_10}.

The receiver must solve three system of linear equation to find the transmission coefficients. We use the ALS algorithm. The ALS algorithm works in the certain order:
The ALS algorithm is very unstable to the ill-conditioned matrices. Possible to add regularization parameters for the algorithm using the Levenberg-Marquardt based improvement.

\subsubsection{rank analysis}\label{part:RA}
 Consider the possible ranks of the generating matrices $\mathbf{\Gamma}_1,\mathbf{\Gamma}_2,\mathbf{\Gamma}_3$ to analyse the decrease size of the task to estimate the all sub-carrier selection coefficients. Two generating matrices consist from two parts corresponding to the symbols and the Khatri-Rao product separation\cite{Book17}.

Consider the the matrix $\mathbf{\Gamma}_1$ \eqref{}. The matrix is based on the Kronecker product of the two matrices. The Kronecker product rank equal to the \eqref{ra_1}, where identity matrix has full rank\cite{Book19}. The rank depends only from the first part of the equation \eqref{ra_2}.

 The upper bound of the rank \eqref{ra_3} is equal to the $min(M_t\cdot F,T)$ due to the size of the resulting matrix. The $T$ is the highest value from the list, so we can neglect them, because it will be higher that any other value in the list\eqref{ra_4}.  The real rank does not depend from the multiplication with rank of the $\mathbf{C}^{[a]'}$ matrix due to the symbol filling of the matrix.  Consider the Khatri-Rao product \cite{Book16}. Row subspace doesn't change from the product due to the problem statement $r\leq F\cdot M_t$, following reveal upper bound of the Khatri-Rao product. Result is shown in the  equation . If the rank of  is full and equal to the $T_s$ the system is bounded with rank $r$, which allow to decrease  rank of the matrix $\mathbf{S}$ from $T_s$ to $r$ and find the $\mathbf{\widehat{A}_A}$.

Next step is $\mathbf{\Gamma}_2$ rank analysis. We doesn't write the overall explanation of the analysis and write only the upper bounds of the rank for each matrix. 

\section{Semi-blind receiver for MIMO}\label{part:SMMIMO}
In the previous subsection was solved task with the transmission coefficient search for the case if the number of the $M_t$ antennas is equal to the $T_s$. Due to high size of transmission blocks, the the $M_t$ is not obligatory equal to the $T_s$. To increase efficiency of the system we can implement the semi-blind receiver from the SISO approach to the MIMO case, with known symbols. Due to ths same PARATUCK2 third order model the algorithm is almost the same. The semi-blind task statement will be written in the same form as in the section  . The modulation matrices will be changed due to the new structure of the signal . Also the equations is changed due to the transposed version of the  of the vectorized symbol matrix and hand vectorized received matrix $\mathbf{Y}$. The all other equations is the same as in the joint semi-blind receiver it the SISO case. The start block of the known symbols increase from the one column, to the $M_t$ columns, for each transmit antenna, and other symbols in the transmission block are filled with data symbols.

\section{The system model for frequency selective channel}\label{part:SMFSC}
In this section we assume the GFDM MIMO system with frequency selective channel and with additive white Gaussian noise at the receiver \eqref{m_fsc_1}\cite{Book29} Explained model is written in the two forms: matrix and vectorized forms.We assume the zero guard interval before transmission block, and zero interference between different blocks. The main equation of the MIMO system with frequency selective channel is written as following \cite{Book53} equation \eqref{m_fsc_1}.

Where the $\mathbf{X}_{:,k}$ is the $k$-th column of the transmitted data block. The size of each column equal of the transmitted antenna size. The $\mathcal{H}$ is the third dimensional tensor and $\mathbf{H}_{:,:,i}$ is $i$-th slice in the third dimension from the tensor $\mathcal{H}$. The $\mathbf{N}_{:,k}$ is the column of the noise matrix We assume that noise has Gaussian distribution, is uncorrelated and has variance $\delta^2$.
This model consider the MIMO case and doesn't include any modulation of the data. The GFDM system also has modulation of the symbols which we should consider in the investigation. We can write modulated data in the two forms, as the transmitted matrix\eqref{m_fsc_2} and as the vectorized notation\eqref{m_fsc_3}. The model is the same as the GFDM MIMO flat frequency model but with changes in the matrix $\mathbf{A}$.

Where the symbol matrix $\mathbf{S}$ is the same as in the GFDM MIMO flat frequency model and constructed from the third dimensional symbol tensor  $\mathcal{S}\in\compl^{F\times T_s\times M_t}$. Also it should be noted, we consider case with $T_sF=T$

The $C^{[b]}$,$\mathbf{C}^{[a]'}$ matrices, $\mathbf{b}$ array stay the same as in the GFDM MIMO flat frequency model .The connection array  stay the same as in the previous model.

The $\mathbf{A}$ matrix changes in comparison with the flat-frequency model. In the original case $\mathbf{A}$ showed connection between receive and transmit antenna. Now it should show transmitted signal. Because in the model we separate each sub-carrier, the $\mathbf{A}$ must sum up all the sub-carriers for each transmit antenna. The following $\mathbf{A}$ from \eqref{m_fsc_9} will sum the sub-carriers for the transmit antennas.

The defined two models allow consider two equation of the frequency selective channels: in the classical \eqref{m_fsc_10} MIMO form and in the vectorized form. The classical form \eqref{m_fsc_10} is written as the sum of different matrix products. The classical form doesn't used in the channel estimation algorithm. The vectorized form \eqref{m_fsc_11} allow to write vectorized received data as the product of the modulation matrix and block Toeplitz convolution matrix. The structure of the matrix $\mathbf{H}$ is presented in the \eqref{m_fsc_12}. where the $\mathbf{H_{x,y}}$ is the convolution matrix for the data which is transmitted from the $y$-th transmit antenna and received with $x$-th receive antenna \eqref{m_fsc_13}. We assume that receiver knows the length of the channel and it equal to the $L+1$ time samples.

We can rewrite equation \eqref{m_fsc_13} with respect to the vectorized channel matrix. We can write the channel values in the following form \eqref{m_fsc_14}

The equation allow to write vectorized received data with respect to the array $\mathbf{h}$.The matrix $\mathbf{D}_{mod}$ is constructed in the special way, similar to the simple convolution matrix. To show the construction process we define the transmitted signal \eqref{m_fsc_16}. The construction is similar to the convolution matrix but for the matrix instead array. After convolution construction the Kronecker product is used to extend matrix for all receive antennas.

\subsubsection{Least squares solution for the channel estimation}
The equation with respect to the transmitted data allow us to construct the optimization problem for the channel vector in the second norm sense\eqref{m_fsc_17}.
To find the minimal residual point we should find the partial derivative of the function and equate it to zero. Because of the complex valued function we used the Wirtinger calculus to find the  partial derivative with respect to the $\mathbf{h}^*$ \eqref{m_fsc_18}. The residual function is convex, and there is only point where partial derivative equal to the zero and this point is global minimum\cite{Book54}. We must solve the system of linear equations to find the solution of equation \eqref{m_fsc_19}. The least squares solution become \eqref{m_fsc_20}.

\subsection{The semi-blind receiver for channel estimation}
The receiver may find the unknown symbol with unknown channel matrix with additional constraints. The approach of the semi-blind receiver must base on the ALS algorithm, because of the different residual functions for each of the unknown array.The main equations for the solution become two residual functions: for the symbols\eqref{m_ce_1} and for the channel coefficients\eqref{m_ce_2}. 

We considered three approaches of the semi-blind receiver: 
\begin{itemize}
\item Simplified semi-blind receiver
\item Normal semi-blind receiver.
\item Newton based semi-blind receiver.
\end{itemize}
The both of them is explained below. Also possible to use the Newton algorithm based approach as in the SISO model here, but to simplify the model and show performance of the both algorithm in the simulation was shown only ALS based algorithm.
\subsection{Simplified semi-blind receiver}
The simplified semi-blind receiver is constructed as the ALS algorithm with typical implementation. The additional information from the known symbols is also renewed in the iteration but ignored in the variables update step. 
The equation for the variables update is shown in \eqref{m_ce_3} and \eqref{m_ce_4} where $\mathbf{S}_{sel}$ is the selection matrix for the unknown symbols matrix elements.

\subsection{ Semi-blind receiver with interference cancellation}
The normal semi-blind receiver use one important property of the received data, the known and unknown symbols possible to separate as the sum between them. The transmitted data can be represented as sum of two components \eqref{m_ce_6} known and unknown. The known and unknown symbol vectors are constructed via multiplication with selection matrix and also possible to change the modulation matrix.

 The receiver can't separate symbols in time domain spatially, but can decrease the interference between known and unknown symbols to archive lower SNR. As the changed objective function for the symbols matrix become \eqref{m_ce_9} function

\subsubsection{Regularization for the semi-blind receiver}
We can add the regularization term in the objective function of semi-blind receiver to make semi-blind receiver more stable in the noise environment. The objective function for the channel estimation will be written as following. The additional regularization term possible to write in the two notations: vectorized and convolution forms.
In the paper \cite{Book53} proposed two algorithms to set the matrix $\mathbf{V}$ for the regularization term:
\begin{itemize}
\item Subspace approach
\item Smoothing least squares approach
\end{itemize}  

\subsection{Newton based semi-blind receiver}
The semi-blind receiver task can be solved via Newton based algorithm. For this objective function should be written in the certain way and we can see that the task can be solved in the similar way as in the SISO case.
 
  The objective function will be written as following\eqref{m_cen_1}. The equation $\mathbf{r}_s$ can be written in the two forms which are equal. The optimal point for the objective function will be point where partial derivative will be equal to the zero \eqref{m_cen_2}.  This point is one because of the  objective function  convexity. We construct the partial derivative with respect to the unknown symbols and channel values. We used the Wirtinger calculus to find the derivative of the complex function \eqref{m_cen_3}.
  The following equations is equated to zero and we solved the system of non-linear equations. To solve the system we used as written above Newton algorithm. 

The Newton algorithm calculate at each iteration solution for the system of linear equations which constructed at the each iteration point \cite{Book68}. We must write the equation of the Jacobian matrix \eqref{m_cen_6} for the  partial derivatives to determine the solved at each iteration equation. To write the Jacobian matrix we use property of the equality of the different notations of the same function in the same point. The resulting Jacobian matrix is written as follows \eqref{m_cen_9}.Unfortunately the Newton algorithm in this statement of the problem is significantly unstable. The typical conditional number of the Jacobian matrix higher than 100. Higher conditional number increase the error in the each iteration. To overcome this problem we propose the simple approach to regularize the problem at each iteration. The explanation of the regularization is written with  the Newton algorithm and is explained below. Regularization allow to make algorithm more stable and converge them even if the Jacobian is very ill-conditioned.
\section{Simulation results}\label{sec:SRMIMO}
In this section we present simulation results for the GFDM MIMO system and algorithms which was reviewed in the chapter \ref{sec:GFDMMIMO}. 

The performance results of the $\mathcal{A}$ search algorithm are obtained through simulation. The parameters of the system are tabulated in Table \ref{tab:m_table_1}.   
The GFDM system is simulated in the AWGN channel without coding. QPSK modulation scheme is used. The number of sub-carriers was $F=16$, samples for each symbol was $T/T_s=F$ The block size was $T_s=5$. The root-raised cosine filter was used with  roll-off factor $\alpha=0.5$. The selection coefficients was chosen mode two product between identity matrix in the first and third dimension with random integer values array.   The results of the GFDM performance is presented in the two plots, at first is shown the reconstruction error for the $\mathbf{A}$ fig. \ref{fig:fs_m_1}, at second is shown the SER in comparison with the case  if $\mathbf{a}=\mathbf{1}$ with the zero-forced receiver fig. \ref{fig:fs_m_2}. 

The GFDM system semi-blind receiver in the MIMO system is simulated in the frequency selective channel without coding. QPSK modulation scheme is used. The channel was generated in the two steps. At the first step was generated flat frequency channel matrix $\mathbf{H}$ with zero-mean circularly symmetric complex Gaussian random values. At the second step for each values of the matrix $\mathbf{H}$ generated the tensor $\mathcal{H}$ with time filter coefficients for each transmission link. At the end each slice of the tensor $\mathcal{H}$ is element-wise multiplied with $\mathbf{H}$. The number of the receive and transmit antennas was $M_r=M_t=3$.  The number of sub-carriers was $F=8$, samples for each symbol was $T/T_s=F$ The block size was $T_s=3$. The root-raised cosine filter was used with roll-of-factor $\alpha=1$. The length of the channel equal to $L+1=4$. The experiment setup is presented in the table \ref{tab:m_table3}. The results of the GFDM performance is presented in the two plots, at first \eqref{fig:fs_m_3} is shown the SER in comparison with the case  if receiver know the original channel and use the Zero-Forced receiver and with case for different number of unknown symbols in one transmission block. At second figure \eqref{fig:fs_m_4} is shown the reconstruction error for the channel values \eqref{sim:eq_m_1}.
The GFDM system Newton based semi-blind receiver in the MIMO system is simulated in the frequency selective channel without coding. QPSK modulation scheme is used. The channel was generated in the two steps. At the first step was generated flat frequency channel matrix $\mathbf{H}$ with zero-mean circularly symmetric complex Gaussian random values. At the second step for each values of the matrix $\mathbf{H}$ generated the tensor $\mathcal{H}$ with time filter coefficients for each transmission link. At the end each slice of the tensor $\mathcal{H}$ is element-wise multiplied with $\mathbf{H}$. The number of the receive and transmit antennas was $M_r=M_t=3$.  The number of sub-carriers was $F=8$, samples for each symbol was $T/T_s=F$ The block size was $T_s=3$. The root-raised cosine filter was used with roll-of-factor $\alpha=1$. The length of the channel equal to $L+1=4$. The experiment setup is presented in the table \ref{tab:m_table3}. The results of the GFDM performance is presented in the two plots, at first \eqref{fig:fs_m_5} is shown the SER in comparison with the case  if receiver know the original channel and use the Zero-Forced receiver and with case for different number of unknown symbols in one transmission block. At second figure \eqref{fig:fs_m_6} is shown the reconstruction error for the channel values. The reconstruction error is defined in the \eqref{sim:eq_m_2}

\section{Conclusion}\label{sec:CMIMO}
The first experiment show that channel orthogonalization also can be done in the another manner with orthogonalization which consider all sub-carriers and channel matrix for the each sub-carrier.

As conclusion for the second experiment with frequency coefficient selection approach in the MIMO case we can see that algorithm decrease SER for 5-6 dB in comparison with system which know the frequency coefficients fig. \ref{fig:fs_m_2}. Decrease in performance come from the simple coefficient selection way in algorithm, where coefficient was chosen as $1$ if the resulting $abs(a_1)$ value is higher that 0.5. There are better solutions to increase performance of the algorithm with the predefined information at the receiver.

There are many frequency coefficients estimation approaches and we compares the separately, but all of them has the same result as we can see from the fig. \ref{fig:fs_m_3}. Even the solution based on the approximated SISO model show the same result. This behaviour come from the additional channel estimation with limit which decrease influence of the noise but also decrease the performance of the more powerful algorithms. Also it should be noted that the joint MIMO algorithm allow to also the symbol estimation, and other algorithms can't estimate the symbols in the received data. The difference between explained algorithms is significant due to the assumptions which was in each of the approaches to estimate $\mathbf{A}$.
 
The main assumption of the SISO semi-blind based approach is error spread between transmit antennas. We ignored unknown values because unknown symbols at the receiver doesn't used. Repetition of the algorithm for the each transmit antenna allow to estimate coefficients for the sub-carrier selection matrix. Assumption is not applied for the time slots due to the overlapping and the receiver estimates the coefficient approximately. The error depends from the difference between sub-carrier selection coefficients in the each transmit antenna due to the time slot mixing in the GFDM systems. The SISO semi-blind based approach is the most expensive algorithm in comparison with others. The receiver can estimate parallel solution for the each sub-carrier selection coefficient corresponding to the each transmit antenna. The parallel solution decreases iteration time of the approach.

The algorithm based on the least squares solution achieve the better results. The approach doesn't have assumptions in the mathematical model  for the GFDM system. But the least squares solution is very sensitive to the condition number of the matrix. The least squares solution is expensive too.

The decomposition based approach of the unfolding multiplication with the matrix decrease size of the problem.  There is trade of between complexity of the task and calculation of the optimization process. Each iteration contain the matrix inversion inside. The matrix to invert change every iteration change from iteration to iteration.
We can see that result of the all three algorithms is the same. There is significant difference for the decomposition based approach if we use the mean value metric. because we can choose bas initial point. Therefore we used the median value metric.

Next experiment is the semi-blind receiver channel and the symbol estimation over frequency selective Rayleigh fading channel. The experiment was done for the all possible number of known symbols which have solution for the problem statement. The result for the SER if shown in the fig. \ref{fig:fs_m_3}. As we can see here, there is one optimal point in symbol estimation. The less unknown symbols in the block increase the SNR, whether the higher number of unknown symbols decrease performance of the approach. The solution where number of unknown symbols equal to $8$ show the closest result to the original GFDM system with the same block size, known channel and zero forced receiver. The channel reconstruction error has the similar behaviour and show the optimal point in the same number of unknown symbols. Important fact is that reconstruction error for the maximal and minimal number of unknown have different behaviour in comparison with the SISO model. But this behaviour also show the convexity of the performance function for the number of unknown symbols as variable. The different block size should have different optimal value of the unknown symbols, but the average performance must be the worse from the known channel system due to the higher energy per bit with same performance. 
The last experiment is the semi-blind receiver channel and the symbol estimation over frequency selective Rayleigh fading channel  with the Newton based approach. The experiment was done for the all possible number of known symbols which have solution for the problem statement. The result for the SER if shown in the fig. \ref{fig:fs_m_5}. 



As we can see here, there is one optimal point in symbol estimation. The less unknown symbols in the block increase the SNR, whether the higher number of unknown symbols decrease performance of the approach. The solution where number of unknown symbols equal to $8$ show the closest result to the original GFDM system with the same block size, known channel and zero forced receiver. The channel reconstruction error has the similar behaviour and show the optimal point in the same number of unknown symbols. Important fact is that reconstruction error for the maximal and minimal number of unknown have different behaviour in comparison with the SISO model. But this behaviour also show the convexity of the performance function for the number of unknown symbols as variable. The different block size should have different optimal value of the unknown symbols, but the average performance must be the worse from the known channel system due to the higher energy per bit with same performance. 

%
\cleardoublepage
\ihead[]{Acronyms}
\printnomenclature


%

\cleardoublepage
\ihead[]{Bibliography}
\bibliographystyle{alphadin}
\bibliography{art1} 
%\appendix

\end{document}

